\newcommand{\PaperIntroduction}{%
\section{Introduction}
\label{sec:introduction}

Generative retrieval (GR) learns to produce a document identifier (DocID) from
a query rather than searching stored document embeddings
\citep{Tay2022TransformerMA,Li2024FromMT}. Among the systems considered here,
IRGen provides the closest structural precedent for a non-text query: a query
image is encoded and a valid database-image identifier is generated
autoregressively
\citep{Zhang2024IRGen}. Its task is same-modality image retrieval, with
relevance defined by shared dataset class labels. GENIUS broadens GR to
text, image, and image--text query content, but each query is conditioned on a
natural-language task instruction and embedded through CLIP-based
vision--language alignment \citep{Kim2025GENIUSAG}. These results establish
non-text-query GR in vision, but not direct spoken-query identifier generation.

The QA retrieval setting introduces a different relevance relation. A spoken
question is paired with a passage because the passage contains the answer, not
because the two recordings are acoustically similar. They may differ in
speaker, duration, wording, and recording conditions. The resulting mapping is
asymmetric: question \(\rightarrow\) answer-containing passage. Spoken retrieval
usually handles this abstraction through ASR followed by text retrieval, or
through dense representations searched in an external vector index
\citep{Wang2023AVATARRV,Lin2024SpeechDPRES}. The former inherits transcription
errors and imposes a text interface; the latter retains embedding search. We
therefore ask: \emph{can a generative retriever map a spoken question directly
to an answer-containing spoken passage when relevance is not defined by acoustic
similarity and no intermediate transcript is available?}

Discrete speech units make this question compatible with sequence models by
turning audio into token sequences without transcribing it
\citep{Hsu2021HuBERTSS,lakhotia-etal-2021-generative,lin22c_interspeech}. The
retriever generates a DocID without querying an external document-embedding
index or performing ANN search. A prefix trie restricts decoding to valid
DocIDs; it is an output constraint rather than a dense retrieval index.

We study \textbf{SpeechGR}, which quantizes spoken questions and passages into
discrete acoustic units and trains a Flan-T5 encoder--decoder to generate the
identifier of an answer-containing passage. The model learns jointly from
passage-to-DocID indexing examples and spoken-query-to-DocID retrieval examples,
and its model inputs consume neither query/passage transcripts nor a textual
query instruction during SpeechGR training or inference. Flan-T5 nevertheless
contributes text pretraining. On the English SLUE-SQA-5 benchmark
\citep{Shon2022SLUEPA},
SpeechGR obtains 5.81 Hit@1 and 18.452 Hit@20 over the 15,883 passages in the
current data export. Speech-unit pretraining provides the strongest positive
experimental pattern, while BPE compression, pairwise ranking loss, and a
windowed Q-Former do not improve their reported references. Gold-text and prior
speech retrievers are reported only as contextual evidence because their
inputs, supervision, mechanisms, or candidate pools differ.

This work makes three contributions:
\begin{itemize}
    \item We study spoken-question generative passage retrieval, in which a
    spoken question generates the identifier of its linked answer-containing
    passage; relevance labels are not based on acoustic similarity.
    \item We develop SpeechGR, which maps discrete speech units to valid
    passage identifiers without consuming task transcripts or textual query
    instructions and retrieves without dense ANN search.
    \item We establish feasibility on a 15,883-passage SLUE-SQA-5 collection
    and report scoped evidence on acoustic units, speech pretraining,
    identifier design, and compression and ranking variants that did not help.
\end{itemize}
}

\newcommand{\PaperRelatedWork}{%
\section{Related Work}
\label{sec:related-work}

\paragraph{Generative retrieval.}
In the standard formulation used here, GR replaces nearest-neighbor search
with autoregressive DocID generation. DSI established the parameterized-index
formulation \citep{Tay2022TransformerMA}; later work studies identifier design,
query generation, and retrieval training
\citep{Zhuang2022BridgingTG,Li2024FromMT,Kuo2024ASO}. Valid-identifier tries
commonly constrain decoding \citep{DeCao2020AutoregressiveER}. Cross-modal GR
generates identifiers for image, audio, and video targets from text queries
\citep{li-etal-2024-generative,fang-etal-2025-cart,
fang25c_interspeech}. IRGen instead maps a raw image to codes for neighboring
images, learns retrieval-specific visual identifiers, and evaluates up to
million-scale collections \citep{Zhang2024IRGen}. It is structurally analogous
to SpeechGR, but evaluates same-modality image retrieval using shared class
labels.
GENIUS accepts visual query content, yet also supplies a textual task
instruction and uses a CLIP-aligned multimodal encoder
\citep{Kim2025GENIUSAG}. SpeechGR instead studies speech-only query content and
an asymmetric question-to-passage relation; its task inputs include neither a
transcript nor a natural-language retrieval instruction.

\paragraph{Spoken retrieval.}
ASR-based pipelines expose downstream retrieval to transcription errors;
AVATAR explicitly studies robustness to such noise
\citep{Wang2023AVATARRV}. SpeechDPR instead learns speech representations for a
dense retriever and distills semantics from an unsupervised ASR model and a
text retriever \citep{Baevski2021UnsupervisedSR,Lin2024SpeechDPRES}. SpeechGR
complements these systems with an ASR-free mapping from a spoken question to an
answer-containing passage DocID rather than using external dense passage
search. Recent dense retrievers report substantially higher SLUE-SQA-5 recall
while importing
strong text-side structure. CLSR optimizes ASR objectives and maps speech to
text-like representations consumed by a frozen BGE encoder
\citep{hu2026clsr}. WavRAG adapts an audio-language model on roughly 1.5
million mixed-modality retrieval examples, including text retrieval corpora
and TTS-derived speech \citep{chen-etal-2025-wavrag}. Their reported candidate
pools are not specified well enough to match SpeechGR's evaluation, but their
training signals clarify that transcription-free inference does not imply
text-independent retrieval learning.

\paragraph{Broader spoken-query benchmarks.}
The Massive Sound Embedding Benchmark (MSEB) introduces Simple Voice Questions
(SVQ), with more than 170k short spoken queries across 26 locales, 17
languages, and four recording conditions. It covers document, passage, and span
retrieval in both in-language and cross-language settings
\citep{Heigold2025MSEB}. Its recordings are readings of XTREME-UP text prompts
and query a textual Wikipedia index. SVQ therefore broadens multilingual and
acoustic-condition coverage, but remains speech-to-text retrieval rather than
speech-to-speech generative DocID retrieval.

\paragraph{A text-derived speech benchmark.}
SLUE-SQA-5 uses natural recordings, but derives questions from five text QA
datasets and pairs passages through transcript-side matching and retrieval
\citep{Shon2022SLUEPA}. Models also rely predominantly on lexical information
when lexical and prosodic cues coexist \citep{chi-etal-2025-role}. We therefore
treat it as a natural-speech realization of a lexically defined retrieval
problem, not a test that isolates speech-specific semantics.

\paragraph{Discrete speech units.}
Self-supervised speech encoders such as HuBERT can be clustered into discrete
units \citep{Hsu2021HuBERTSS}. Such units support generative spoken-language
modeling and downstream spoken understanding without requiring transcripts as
model inputs \citep{lakhotia-etal-2021-generative,lin22c_interspeech,
Shon2024DiscreteSLUAL}. We use them as the input vocabulary of a generative
retriever and examine how layer choice, codebook size, and sequence compression
affect retrieval.
}

\newcommand{\PaperLimitations}{%
\section{Limitations}
\label{sec:limitations}

The evidence is limited to English SLUE-SQA-5 and its linked Spoken Wikipedia
collection, so it does not establish multilingual performance, effectiveness
for unwritten languages, or scaling behavior on substantially larger
collections. The results also do not establish semantic reasoning. SLUE-SQA-5
is derived from text QA data and
transcript-based passage selection, so lexical correspondence may account for
part of the observed retrieval. Controlled paraphrase, low-overlap, and
acoustic-distractor evaluations would be needed to separate semantic mapping
from lexical pattern matching. SpeechGR is task-transcript-free rather than
text-independent: Flan-T5 is text-pretrained and the identifiers are
metadata-derived. IRGen and GENIUS are structural, not numerical, comparators:
class-based image retrieval, instruction-conditioned multimodal search, and
question-to-passage retrieval use different relevance labels and scales. The comparison with
SpeechDPR is neither supervision- nor candidate-pool-matched: SpeechGR uses
15,883 paired passages, whereas SpeechDPR reports retrieval over roughly 39k
Spoken Wikipedia segments. The CLSR and WavRAG candidate pools are
under-specified, and all of their values are imported rather than reproduced.
The gold-text BM25 reference uses one fixed tokenizer and parameter setting and
is not a sparse-retrieval sensitivity study.
The thesis does not report repeated runs, confidence intervals, or significance
tests.
The experimental record also does not preserve the per-ablation split and shared
pretraining setting, input-truncation policy, indexing/retrieval sampling
ratio, optimization hyperparameters, or completed-beam scoring needed for
exact reproduction; it is internally inconsistent about whether the
\(k\)-means codebook used 100 or 960 hours of LibriSpeech. Small differences
should therefore be interpreted cautiously. Finally, the experiments do not
evaluate corpus updates or retraining cost when passages change.
}

\newcommand{\PaperConclusion}{%
\section{Conclusion}
\label{sec:conclusion}

SpeechGR transfers the identifier-generation structure used in visual
retrieval to an asymmetric spoken task: a spoken question generates the DocID
of an answer-containing spoken passage. It uses no task transcript, textual query
instruction, or external document-embedding index. Its 5.81 Hit@1 and 18.452
Hit@20 on 15,883 SLUE-SQA-5 passages establish feasibility rather than parity
with text-guided retrieval. Speech-unit pretraining gives the strongest
positive pattern; the tested BPE, ranking-loss, and Q-Former variants do not.
Because the benchmark is text-derived, these results show learnable
question-to-passage associations but do not isolate semantic reasoning from
lexical matching. That distinction is the central challenge for future spoken
generative retrieval.
}
